\documentclass[a4paper,12pt]{article}

% --- Pakiety ---
\usepackage[polish]{babel}
\usepackage[utf8]{inputenc}
\usepackage[T1]{fontenc}
\usepackage[top=2.5cm,bottom=2.5cm,left=3cm,right=3cm]{geometry}
\usepackage{amsmath}
\usepackage{graphicx}
\usepackage{booktabs}
\usepackage[colorlinks=true, allcolors=blue]{hyperref}
\usepackage{amssymb}
\usepackage{listings}
\usepackage{titlesec}

\lstset{
    basicstyle=\ttfamily\small,
    breaklines=true,
    columns=fullflexible
}

\titlespacing*{\subsection}{0pt}{12pt}{6pt}
\setlength{\parindent}{15pt}
\setlength{\parskip}{6pt}

% --- Tytuł ---
\title{
    \vspace{1.5cm}
    \textbf{PODSTAWY TEORII INFORMACJI}\\
    \vspace{1.5cm}
    Laboratorium 3\\
    Syntaktyczna analiza danych\\
    \vspace{1.5cm}
    \huge \textbf{Ekstrakcja i redukcja informacji w danych meteorologicznych}\\
    \vspace{1cm}
}

\author{
    Maciek Chotkowski \\ Martsin Lazouski \\ Nadia Ziecina \\
    Michał Kasznia \\ Kasiński Julian \\ Volodymyr Kamenchuk \\ Usevalad Nikalaichuk 
}

\date{\today}

\begin{document}

\maketitle
\newpage

\tableofcontents
\newpage

%--------------------------------------------------

\section{Wprowadzenie}

W poprzednim etapie projektu konieczne było sprowadzenie analizowanego problemu do uproszczonej postaci, umożliwiającej dalsze przetwarzanie danych. Ograniczono się do wybranych parametrów meteorologicznych oraz ich reprezentacji w postaci ciągów liczbowych.

Dopiero w niniejszym etapie wprowadzamy konkretne rozwiązanie obliczeniowe. Oznacza to przejście od opisu problemu do jego realizacji w postaci narzędzia przetwarzającego dane.

Celem ćwiczenia jest analiza syntaktyczna danych, polegająca na ich przetwarzaniu i ocenie przy użyciu formalnych kryteriów obliczeniowych.

%--------------------------------------------------

\section{Etap 1: Narzędzie pozyskiwania danych}

Pierwszym etapem opracowanego systemu jest moduł automatycznego pobierania
danych meteorologicznych i zapisania ich w prostej strukturze plików.
Na tym etapie nie tworzymy jeszcze rankingu ani jednego wspólnego pliku
porównawczego. Celem jest jedynie zebranie danych i rozdzielenie ich według
źródła oraz typu wielkości fizycznej.

Każde źródło danych otrzymuje własny katalog. W katalogu znajdują się trzy
pliki CSV:
\begin{verbatim}
nazwa_zrodla/
    temp.csv
    opady.csv
    wiatr.csv
\end{verbatim}

Taki zapis jest zgodny z poziomem syntaktycznym analizy, ponieważ każdy plik
zawiera prosty ciąg liczbowy jednej cechy. Temperatura, opady i wiatr nie są
jeszcze interpretowane semantycznie. Są traktowane jako sygnały liczbowe,
które można później dopasować po czasie i porównać z innymi źródłami.

Wejściem do modułu są:
\begin{itemize}
    \item współrzędne geograficzne,
    \item horyzont prognozy,
    \item nazwa źródła danych.
\end{itemize}

Wyjściem są trzy uporządkowane sygnały:
\[
T = \{T_1, T_2, \dots, T_n\}, \quad
P = \{P_1, P_2, \dots, P_n\}, \quad
W = \{W_1, W_2, \dots, W_n\}
\]

Dane są zapisywane osobno dla każdego źródła i przekazywane do dalszych
etapów analizy.
\newpage
\subsection{Implementacja narzędzia}

Poniższy fragment przedstawia podstawową wersję programu pozyskującego dane
z dwóch źródeł: Open-Meteo oraz yr.no. Program pobiera z odpowiedzi API tylko
pola potrzebne w projekcie: czas, temperaturę, opady i prędkość wiatru.
Następnie zapisuje każdą wielkość do osobnego pliku w katalogu danego
źródła.

Wynikiem działania kodu nie jest jeden plik zbiorczy. Wynikiem jest katalog
\texttt{open\_meteo} oraz katalog \texttt{yr\_no}. W każdym z nich znajdują
się pliki \texttt{temp.csv}, \texttt{opady.csv} oraz \texttt{wiatr.csv}.
Dzięki temu po etapie zbierania danych wiadomo, z którego źródła pochodzi
każdy sygnał i jakiego parametru dotyczy.

\begin{lstlisting}[language=Python]
import requests, pandas as pd
from pathlib import Path

LAT, LON = 52.23, 21.01

om = requests.get(
    "https://api.open-meteo.com/v1/forecast",
    params={
        "latitude": LAT,
        "longitude": LON,
        "hourly": "temperature_2m,precipitation,wind_speed_10m",
        "forecast_days": 3
    }
).json()

df_om = pd.DataFrame({
    "time": om["hourly"]["time"],
    "temp": om["hourly"]["temperature_2m"],
    "opady": om["hourly"]["precipitation"],
    "wiatr": om["hourly"]["wind_speed_10m"]
})

yr = requests.get(
    "https://api.met.no/weatherapi/locationforecast/2.0/compact",
    headers={"User-Agent": "pti-lab"},
    params={"lat": LAT, "lon": LON}
).json()

df_yr = pd.DataFrame([{
    "time": item["time"][:16],
    "temp": item["data"]["instant"]["details"]["air_temperature"],
    "opady": item["data"].get("next_1_hours", {}).get("details", {}).get("precipitation_amount", 0),
    "wiatr": item["data"]["instant"]["details"]["wind_speed"]
} for item in yr["properties"]["timeseries"]])

def zapisz_zrodlo(nazwa, df):
    out = Path(nazwa)
    out.mkdir(exist_ok=True)
    df[["time", "temp"]].to_csv(out / "temp.csv", index=False)
    df[["time", "opady"]].to_csv(out / "opady.csv", index=False)
    df[["time", "wiatr"]].to_csv(out / "wiatr.csv", index=False)

zapisz_zrodlo("open_meteo", df_om)
zapisz_zrodlo("yr_no", df_yr)
\end{lstlisting}

\subsection{Zakres i postać pobieranych danych}

W ramach narzędzia pobierane są prognozowane dane meteorologiczne
dla lokalizacji Warszawa (52.23N, 21.01E).

Uwzględniono następujące parametry:
\begin{itemize}
    \item temperatura powietrza (°C),
    \item opady (mm),
    \item prędkość wiatru (m/s).
\end{itemize}

Dane pobierane są z API w formacie JSON, który ma strukturę
hierarchiczną (zagnieżdżone pola i listy). Poszczególne wartości
odczytywane są z odpowiednich pól, np.:

\begin{verbatim}
Open-Meteo: hourly -> temperature_2m
yr.no: properties -> timeseries -> instant -> details
\end{verbatim}

Po ekstrakcji dane przekształcane są do postaci tabelarycznej
oraz wektorowej:

\[
T(t), \quad P(t), \quad W(t)
\]

gdzie:
\begin{itemize}
    \item $T(t)$ – temperatura,
    \item $P(t)$ – opady,
    \item $W(t)$ – prędkość wiatru.
\end{itemize}

Ostatecznie dane nie są jeszcze łączone w jeden plik. Każdy parametr jest
zapisywany osobno w katalogu źródła:
\begin{verbatim}
open_meteo/temp.csv
open_meteo/opady.csv
open_meteo/wiatr.csv
yr_no/temp.csv
yr_no/opady.csv
yr_no/wiatr.csv
\end{verbatim}
Analogicznie można zapisać kolejne źródła, np. \texttt{graphcast}
albo \texttt{stacja\_pomiarowa}. Dopiero w następnym
etapie pliki te są dopasowywane po czasie i porównywane.

\subsection{Rozszerzenie: pobieranie i wektoryzacja danych z modeli GraphCast, GenCast i WeatherNext}

Oprócz klasycznych serwisów pogodowych można wykorzystać prognozy tworzone
przez modele uczenia maszynowego: GraphCast, GenCast oraz WeatherNext.
W raporcie traktujemy je jako dodatkowe źródła danych, ale po pobraniu
muszą one zostać sprowadzone do dokładnie takiej samej postaci jak dane
z Open-Meteo i yr.no. Oznacza to, że końcowym wynikiem nie jest duży plik
z globalną prognozą, tylko proste wektory wartości dla wybranej lokalizacji:
\[
temp = [temp_1, temp_2, \dots, temp_n]
\]
\[
precip = [precip_1, precip_2, \dots, precip_n]
\]
\[
wind = [wind_1, wind_2, \dots, wind_n]
\]

GraphCast jest modelem deterministycznym, więc dla jednego czasu startu
zwraca jedną prognozę. GenCast jest modelem zespołowym, dlatego może zwrócić
wiele wariantów tej samej prognozy. WeatherNext udostępnia natomiast gotowe
operacyjne prognozy w usługach Google, między innymi przez BigQuery, Earth
Engine oraz pliki Zarr w Google Cloud Storage. Z punktu widzenia naszego
kodu wszystkie te źródła powinny jednak przejść przez ten sam etap
standaryzacji.

W praktyce przyjmujemy następujący układ plików:
\begin{itemize}
    \item \texttt{open\_meteo/temp.csv}, \texttt{open\_meteo/opady.csv}, \texttt{open\_meteo/wiatr.csv} -- dane z Open-Meteo,
    \item \texttt{yr\_no/temp.csv}, \texttt{yr\_no/opady.csv}, \texttt{yr\_no/wiatr.csv} -- dane z yr.no,
    \item \texttt{ai\_graphcast\_vectors.csv} -- wektory wyciągnięte z prognozy GraphCast,
    \item \texttt{ai\_gencast\_vectors.csv} -- wektory wyciągnięte z prognozy GenCast,
    \item \texttt{ai\_weathernext\_vectors.csv} -- wektory wyciągnięte z prognozy WeatherNext,
    \item \texttt{weather\_ai\_vectors.csv} -- wspólny plik, w którym znajdują się wszystkie źródła AI,
    \item \texttt{weather\_all\_vectors.csv} -- opcjonalny plik łączący dane klasyczne i dane AI.
\end{itemize}

Najważniejsze jest to, że każdy plik z danymi AI ma tę samą strukturę
kolumn:
\begin{verbatim}
time, source, member, temp_c, precip_mm, wind_ms, pressure_hpa
\end{verbatim}

Kolumna \texttt{source} przechowuje nazwę źródła, np. \texttt{GraphCast},
\texttt{GenCast} albo \texttt{WeatherNext}. Kolumna \texttt{member} jest
potrzebna głównie dla GenCast, ponieważ model zespołowy może wygenerować
kilka wariantów prognozy. Dla GraphCast i pojedynczej średniej prognozy
WeatherNext można wpisać wartość \texttt{0}. Dzięki temu w dalszej analizie
nie tworzymy osobnego kodu dla każdego modelu, tylko filtrujemy dane po
kolumnie \texttt{source}.

\subsubsection{Przekształcenie danych globalnych na proste wektory}

Modele AI nie zwracają od razu pojedynczej tabeli dla Warszawy. Ich dane
mają postać pól meteorologicznych na siatce geograficznej. Dlatego przed
zapisaniem CSV wykonujemy kilka stałych kroków:
\begin{enumerate}
    \item wybieramy czas startu prognozy oraz horyzont, np. 72 godziny,
    \item wybieramy najbliższy punkt siatki dla Warszawy \((52.23N, 21.01E)\),
    \item odczytujemy potrzebne zmienne meteorologiczne,
    \item przeliczamy jednostki na takie same jak w pozostałych plikach,
    \item z komponentów wiatru wyliczamy jedną prędkość wiatru,
    \item zapisujemy wynik jako zwykłą tabelę CSV.
\end{enumerate}

Najważniejsze konwersje są następujące:
\begin{itemize}
    \item temperatura: z kelwinów na stopnie Celsjusza, czyli \texttt{temp\_c} = \texttt{temp\_k} - 273.15,
    \item opady: z metrów na milimetry, czyli \texttt{precip\_mm} = \texttt{precip\_m} $\cdot$ 1000,
    \item wiatr: z dwóch składowych na prędkość, czyli \texttt{wind\_ms} = $\sqrt{u^2 + v^2}$,
    \item ciśnienie: z paskali na hektopaskale, czyli \texttt{pressure\_hpa} = \texttt{pressure\_pa} / 100.
\end{itemize}

Po takim przekształceniu dane AI są używane w raporcie tak samo jak dane
z klasycznych API. Różnica polega tylko na tym, że źródłem wartości jest
model globalny, a nie endpoint zwracający prognozę od razu dla jednego
miasta.

\subsubsection{Parametry modeli wykorzystywane w projekcie}

\begin{table}[h]
\centering
\small
\begin{tabular}{p{3.0cm}p{4.3cm}p{6.1cm}}
\toprule
Parametr w raporcie & Nazwa w danych AI & Wykorzystanie w kodzie \\
\midrule
Temperatura & \texttt{2m\_temperature} & Wektor \texttt{temp\_c}; porównanie z pomiarami rzeczywistymi i obliczanie RMSE. \\
Opady & \texttt{total\_precipitation\_6hr} lub \texttt{total\_precipitation} & Wektor \texttt{precip\_mm}; analiza sum opadów oraz błędu prognozy. \\
Wiatr & \texttt{10m\_u\_component\_of\_wind}, \texttt{10m\_v\_component\_of\_wind} & Wektor \texttt{wind\_ms}; prędkość liczona z dwóch składowych kierunkowych. \\
Ciśnienie & \texttt{mean\_sea\_level\_pressure} & Wektor \texttt{pressure\_hpa}; parametr dodatkowy, przydatny przy rozszerzeniu analizy. \\
Wilgotność & \texttt{specific\_humidity} & Parametr opcjonalny; może pomóc wyjaśniać błędy temperatury i opadów. \\
Geopotencjał & \texttt{geopotential} & Parametr opcjonalny; opisuje strukturę atmosfery na poziomach ciśnienia. \\
\bottomrule
\end{tabular}
\caption{Parametry modeli AI i sposób ich użycia w projekcie}
\end{table}

GraphCast i GenCast udostępniają modele oraz przykładowe dane przez
repozytorium Google DeepMind GraphCast\footnote{\url{https://github.com/google-deepmind/graphcast}}.
WeatherNext udostępnia prognozy jako gotowe zbiory danych, między innymi
w BigQuery, Earth Engine oraz Google Cloud Storage/Zarr\footnote{\url{https://developers.google.com/weathernext/guides/access-forecast}}.
W dokumentacji WeatherNext 2 podane są również nazwy zmiennych, rozdzielczość
0.25 stopnia, horyzont do 15 dni oraz dane powierzchniowe takie jak
\texttt{2m\_temperature}, \texttt{10m\_u\_component\_of\_wind},
\texttt{10m\_v\_component\_of\_wind}, \texttt{total\_precipitation\_6hr}
i \texttt{mean\_sea\_level\_pressure}\footnote{\url{https://developers.google.com/weathernext/guides/model-specs-vmg}}.

\subsubsection{Przykładowa implementacja wektoryzacji}

Poniższy kod pokazuje, w jaki sposób można ujednolicić dane z różnych
modeli AI. Zakładamy, że wynik działania GraphCast i GenCast został zapisany
do pliku Zarr lub NetCDF, a WeatherNext jest czytany bezpośrednio z Google
Cloud Storage. Każde źródło przechodzi przez tę samą funkcję
\texttt{vectorize\_ai\_source}, która zwraca tabelę z prostymi wektorami.

Opis procesu znajduje się przed listingiem, dlatego sam kod zawiera wyłącznie
instrukcje programu. Funkcja otwiera zbiór danych, wybiera punkt najbliższy
Warszawie, odczytuje wymagane zmienne, uśrednia warianty zespołowe oraz
zapisuje wynik we wspólnym formacie. Ten etap jest przykładem aproksymacji:
z wielowymiarowego pola atmosferycznego zostaje wybrany lokalny wektor cech
potrzebny do rankingu.

\begin{lstlisting}[language=Python]
import requests
import pandas as pd
from pathlib import Path

LAT = 52.23
LON = 21.01

DATA_START = "2026-05-10"
DATA_END = "2026-05-13"

MODELS = {
    "graphcast": "gfs_graphcast025",
    "aifs": "ecmwf_aifs025_single"

}

for source_name, model_name in MODELS.items():

    url = "https://api.open-meteo.com/v1/forecast"

    params = {
        "latitude": LAT,
        "longitude": LON,
        "start_date": DATA_START,
        "end_date": DATA_END,
        "hourly": (
            "temperature_2m,"
            "precipitation,"
            "wind_speed_10m"
        )
    }

    if model_name is not None:
        params["models"] = model_name

    response = requests.get(
        url,
        params=params
    )

    data = response.json()

    if "hourly" not in data:

        print(f"Błąd dla {source_name}")
        print(data)
        continue

    df = pd.DataFrame({
        "time": data["hourly"]["time"],
        "temp": data["hourly"]["temperature_2m"],
        "opady": data["hourly"]["precipitation"],
        "wiatr": data["hourly"]["wind_speed_10m"]
    })

    out = Path(source_name)
    out.mkdir(exist_ok=True)

    df[["time", "temp"]].to_csv(
        out / "temp.csv",
        index=False
    )

    df[["time", "opady"]].to_csv(
        out / "opady.csv",
        index=False
    )

    df[["time", "wiatr"]].to_csv(
        out / "wiatr.csv",
        index=False
    )

    print(f"Zapisano dane dla: {source_name}")
\end{lstlisting}

Jeżeli chcemy zachować wszystkie warianty GenCast, nie należy wykonywać
uśredniania po wymiarze \texttt{member}. Wtedy każdy członek zespołu jest
zapisany jako osobny zestaw wektorów w tym samym pliku CSV. Daje to możliwość
policzenia nie tylko błędu średniej prognozy, ale również rozrzutu prognoz:
mały rozrzut oznacza dużą zgodność wariantów, a duży rozrzut oznacza większą
niepewność modelu.

\subsubsection{Jak rozszerzyć system o kolejne dane AI}

Rozszerzenie kodu o następny model AI polega na dopisaniu nowego źródła do
słownika \texttt{AI\_SOURCES} i ewentualnym uzupełnieniu list nazw zmiennych
w funkcji \texttt{choose\_var}. Reszta programu nie musi się zmieniać,
ponieważ każdy model zostaje przekształcony do tego samego schematu kolumn.

Najbardziej praktyczne rozszerzenia to:
\begin{itemize}
    \item dodanie kolejnych lokalizacji i zapis kolumn \texttt{lat}, \texttt{lon},
    \item zapis surowych członków zespołu GenCast zamiast samej średniej,
    \item dołączenie wilgotności i ciśnienia do analizy błędów,
    \item interpolacja danych AI do godzin wspólnych z Open-Meteo i yr.no,
    \item automatyczne tworzenie pliku \texttt{weather\_all\_vectors.csv}, który łączy wszystkie źródła.
\end{itemize}

Wpływ na dalszy kod jest więc kontrolowany: zamiast pisać osobne skrypty
dla GraphCast, GenCast i WeatherNext, dopisujemy jedną warstwę
standaryzującą. MATLAB lub Python analizują później tylko pliki CSV z
kolumnami \texttt{time}, \texttt{source}, \texttt{temp\_c},
\texttt{precip\_mm} oraz \texttt{wind\_ms}. Dzięki temu porównanie RMSE,
filtracja średnią ruchomą i tworzenie wykresów działają tak samo dla danych
klasycznych oraz danych AI.

\section{Etap 2: Przetwarzanie i analiza danych}

W drugim etapie przeprowadzono przetwarzanie danych oraz ich analizę syntaktyczną.

Dane zapisane w katalogach źródeł zostały wczytane do środowiska MATLAB,
a następnie połączone po kolumnie \texttt{time}. Dopiero w tym etapie
oddzielne pliki \texttt{temp.csv}, \texttt{opady.csv} i \texttt{wiatr.csv}
są składane w jedną tabelę roboczą. Dzięki temu etap pozyskiwania danych
pozostaje prosty, a etap porównywania może dowolnie wybierać źródła.

Po rozszerzeniu projektu o modele AI analizujemy dwa typy reprezentacji:
katalogi źródeł klasycznych oraz plik \texttt{weather\_ai\_vectors.csv},
czyli dane AI sprowadzone do prostych wektorów temperatury, opadów i wiatru.

%--------------------------------------------------
\subsection{Implementacja w MATLAB}

Kod w MATLAB-ie realizuje etap porządkowania i oceny informacji. Program
wczytuje trzy typy danych: prognozy klasyczne, dane AI oraz dane rzeczywiste.
Następnie wszystkie znaczniki czasu są konwertowane do formatu
\texttt{datetime}, co umożliwia wykonanie złączenia tabel po kolumnie
\texttt{time}. Po takim dopasowaniu porównywane są wartości odnoszące się do
tych samych chwil.

W pierwszej części liczony jest błąd dla klasycznego źródła prognozy.
Następnie program przechodzi po kolejnych modelach AI zapisanych w kolumnie
\texttt{source}. Każde źródło jest oceniane tym samym wzorem RMSE, dzięki
czemu porównanie ma charakter formalny i obliczeniowy. Na końcu tworzone są
wykresy, które pokazują zarówno wartości prognozowane, jak i różnicę między
prognozą a danymi referencyjnymi.

\begin{lstlisting}[language=Matlab]
model = [
    wczytaj_zrodlo('open_meteo', 'Open-Meteo');
    wczytaj_zrodlo('yr_no', 'yr.no')
];
ai = readtable('weather_ai_vectors.csv');
real = readtable('real_data.csv');

model.source = string(model.source);
model.time = datetime(model.time);
ai.source = string(ai.source);
ai.time = datetime(ai.time);
real.time = datetime(real.time);
rmse = @(a,b) sqrt(mean((a - b).^2));

sources = unique(model.source);

for i = 1:numel(sources)
    src = sources(i);
    S = model(model.source == src, :);
    T = innerjoin(S, real, 'Keys', 'time');

    rmse_temp = rmse(T.temp, T.temp_real);
    rmse_wind = rmse(T.wiatr, T.wind_real);
    rmse_opady = rmse(T.opady, T.precip_real);

    fprintf('RMSE temperatura %s: %.2f\n', src, rmse_temp);
    fprintf('RMSE wiatr %s: %.2f\n', src, rmse_wind);
    fprintf('RMSE opady %s: %.2f\n', src, rmse_opady);
end

ai_sources = unique(ai.source);

for i = 1:numel(ai_sources)
    src = ai_sources(i);

    S = ai(ai.source == src & ai.member == 0, :);
    T_ai = innerjoin(S, real, 'Keys', 'time');

    rmse_temp_ai = rmse(T_ai.temp_c, T_ai.temp_real);
    rmse_wind_ai = rmse(T_ai.wind_ms, T_ai.wind_real);
    rmse_precip_ai = rmse(T_ai.precip_mm, T_ai.precip_real);

    fprintf('RMSE temperatura %s: %.2f\n', src, rmse_temp_ai);
    fprintf('RMSE wiatr %s: %.2f\n', src, rmse_wind_ai);
    fprintf('RMSE opady %s: %.2f\n', src, rmse_precip_ai);
end

function T = wczytaj_zrodlo(folder, source)
    temp = readtable(fullfile(folder, 'temp.csv'));
    opady = readtable(fullfile(folder, 'opady.csv'));
    wiatr = readtable(fullfile(folder, 'wiatr.csv'));

    T = innerjoin(temp, opady, 'Keys', 'time');
    T = innerjoin(T, wiatr, 'Keys', 'time');
    T.source = repmat(string(source), height(T), 1);
end

\end{lstlisting}

Dodanie danych AI wpływa więc na kod przede wszystkim przez zmianę warstwy
wejściowej. Nie zmienia się sposób liczenia RMSE ani filtracji, ponieważ
każdy model AI zostaje wcześniej sprowadzony do kolumn \texttt{temp\_c},
\texttt{precip\_mm} i \texttt{wind\_ms}. Zamiast pisać oddzielne obliczenia
dla GraphCast, GenCast i WeatherNext, MATLAB przechodzi po liście wartości
w kolumnie \texttt{source} i wykonuje ten sam zestaw operacji dla każdego
źródła.

\subsection{Analiza z wykorzystaniem danych rzeczywistych}

Dane prognozowane zostały zestawione z rzeczywistymi pomiarami
meteorologicznymi zapisanymi w pliku CSV.

Po dopasowaniu danych względem czasu przeprowadzono analizę
różnic pomiędzy prognozą a rzeczywistością.

Na tej podstawie obliczono błąd RMSE, który stanowi miarę
dokładności prognozy.


%--------------------------------------------------
\subsection{Wyniki}

\begin{table}[h]
\centering
\begin{tabular}{lc}
\toprule
Parametr & RMSE \\
\midrule
Temperatura & X \\
Temperatura (po filtracji) & Y \\
Wiatr & Z \\
Opady & W \\
\bottomrule
\end{tabular}
\caption{Błąd RMSE przed i po filtracji}
\end{table}

%--------------------------------------------------
\subsection{Interpretacja}

Filtracja danych przy użyciu średniej ruchomej prowadzi do redukcji
lokalnych fluktuacji sygnału.

W przypadku temperatury obserwuje się zmniejszenie błędu RMSE
po zastosowaniu filtracji, co wskazuje na obecność szumu w danych.

Dla wiatru oraz opadów różnice pozostają większe,
co wynika z ich bardziej dynamicznego i nieregularnego charakteru.

%--------------------------------------------------
\subsection{Wnioski}

Zastosowane metody przetwarzania pozwalają na poprawę jakości
reprezentacji danych poprzez redukcję szumu.

Jednocześnie wyniki pokazują, że skuteczność metod zależy
od charakteru analizowanej zmiennej.

Temperatura jest najlepiej przewidywalna,
natomiast opady pozostają najbardziej niestabilne.
\end{document}
